% BHU PYQ -- Professional Exam Template
\documentclass[11pt,a4paper]{article}

\usepackage[a4paper,top=2.5cm,bottom=2.5cm,left=2.8cm,right=2.8cm,headheight=14pt]{geometry}
\usepackage{amsmath,amssymb}
\usepackage[T1]{fontenc}
\usepackage[utf8]{inputenc}
\usepackage{lmodern}
\usepackage{microtype}
\usepackage{fancyhdr}
\usepackage{enumitem}
\usepackage{listings}
\usepackage{hyperref}
\usepackage{lastpage}
\usepackage{parskip}

\hypersetup{colorlinks=true,urlcolor=black,linkcolor=black,
  pdftitle={CSB-602: Data and File Structures -- BHU PYQ}}

\pagestyle{fancy}
\fancyhf{}
\renewcommand{\headrulewidth}{0.4pt}
\renewcommand{\footrulewidth}{0.4pt}
\fancyhead[L]{\small   \textit{Banaras Hindu University}}
\fancyhead[R]{\small   \textit{CSB-602: Data and File Structures}}
\fancyfoot[C]{\small Page  \thepage\ of \pageref{LastPage}}


\newlist{parts}{enumerate}{2}
\setlist[parts,1]{label=(\alph*),leftmargin=*,itemsep=4pt,topsep=3pt}
\setlist[parts,2]{label=(\roman*),leftmargin=*,itemsep=2pt,topsep=2pt}

\newcommand{\pts}[1]{\hfill   \textit{[#1]}}


\begin{document}

\begin{center}
  {\large\bfseries BANARAS HINDU UNIVERSITY}\\[2pt]
  {\normalsize B.Sc.\ (Hons.) Semester VI Examination, 2013--14}\\[6pt]
  \rule{0.8\linewidth}{0.5pt}\\[6pt]
  {\large\bfseries Paper: CSB-602 --- Data and File Structures}\\[4pt]
  {\normalsize Time: Three Hours \hfill Maximum Marks: 70}
\end{center}

\rule{\linewidth}{0.4pt}

\smallskip



\noindent   \textit{Instructions: Answer   \textbf{five} questions in all, including Question~1 which is compulsory. Figures on the right-hand side indicate marks.}

\bigskip




\noindent   \textbf{Question 1.} Answer   \textbf{all} parts:\pts{$3 imes6=18$}

\begin{parts}
  \item Define the terms: Abstract Data Types (ADT) and Data Structure.
  \item State the advantages of using postfix notation over infix notation.
  \item Define Double-ended Queue (Deque).
  \item What is a Binary Heap? Distinguish between Max-Heap and Min-Heap.
  \item List the basic operations carried out on a linked list.
  \item Briefly discuss the following terms: Tree, Degree of a node, and Height of a tree.
  \item What do you mean by collision in hashing? Briefly explain any two collision resolution methods.
\end{parts}

\bigskip




\noindent   \textbf{Question 2.}\pts{$2+5+6=13$}

\begin{parts}
  \item Define ADT Array.
  \item What is a sparse matrix? Give a structure for its efficient representation.
  \item Give an algorithm for addition of two sparse matrices stored in the representation you proposed above.
\end{parts}

\bigskip




\noindent   \textbf{Question 3.}\pts{$2+6+5=13$}

\begin{parts}
  \item What are the two important criteria to be satisfied by any recursive function?
  \item Explain the implementation of a stack using an array with algorithms for all operations on a stack.
  \item Write an algorithm to convert an infix expression to a postfix expression, and explain it through an example.
\end{parts}

\bigskip




\noindent   \textbf{Question 4.}\pts{$3+10=13$}

\begin{parts}
  \item What do you mean by order of growth in reference to time complexity of an algorithm?
  \item Explain the quick sort algorithm with an example. Also mention any three strategies for choosing a pivot element.
\end{parts}

\bigskip




\noindent   \textbf{Question 5.}\pts{$4+9=13$}

\begin{parts}
  \item Define full binary tree and complete binary tree with an example for each.
  \item Explain binary tree traversals (in-order, pre-order, post-order) with recursive algorithms.
\end{parts}

\bigskip




\noindent   \textbf{Question 6.}\pts{$7+6=13$}

\begin{parts}
  \item Explain Breadth First Search (BFS) traversal on a graph with an example.
  \item Write an algorithm for adding two polynomials using a linked list.
\end{parts}

\bigskip




\noindent   \textbf{Question 7.} Write notes on any   \textbf{two} of the following:\pts{$6.5 imes2=13$}

\begin{parts}
  \item Hashing
  \item Binary Search Algorithm (Iterative and Recursive)
  \item Sequential and Index-Sequential file organization techniques
  \item External sorting methods with an example
\end{parts}

\vfill
\rule{\linewidth}{0.4pt}

\begin{center}\small
\href{https://bhu-exam.vercel.app/}{Click here to visit our website}\\[4pt]\href{https://me-aryan.vercel.app/}{Click here to visit our contributor}\\[4pt]\href{https://quashan-me.vercel.app/}{Click here to visit our contributor}
\end{center}
\end{document}